\section{Discussion}
\label{sec:discussion}

We have developed an exact polynomial-time simulation framework for a broad class of non-Clifford circuits, with no restriction on the number of logical qubits or non-Clifford gates.
It rests on two equivalent descriptions: third-order phase-polynomial states, specified by affine supports and polynomial phases, and DCP stabilizer states, which extend the stabilizer picture beyond Pauli operators and yield a tableau-based simulator.
We characterize precisely which Pauli measurements preserve this state class and provide efficient tests and update rules supporting intermediate measurements and classical feedforward.

For quantum error correction, a simple condition on the ideal circuit guarantees compatibility under arbitrary inserted DCP faults: each ideal $X$ measurement must either be deterministic or anticommute with a fixed $Z$-type constraint.
This remains sufficient even when noise makes deterministic outcomes random, establishing efficient simulation for protocols involving magic state preparation, non-Pauli checks, and CSS gauge fixing.

Our open-source implementation, \texttt{merlin}~\cite{koch2026merlin}, makes this framework available for circuit-level simulation.
Benchmarks show favorable runtime and memory scaling with the number of logical outputs in Bravyi-Haah distillation, and successful simulation of the distance-5 code switching circuit on which all other tested simulators exhaust memory or fail.
On the other hand, existing methods currently achieve lower per-shot runtimes on the single-output distillation and cultivation examples.

A natural next step is to use this framework to evaluate recent proposals for non-Clifford computation with quantum LDPC codes, including transversal gates, code switching, and high-rate magic state preparation, whenever their circuit implementations satisfy our compatibility conditions.
Combined with suitable decoders, our simulator could enable direct comparisons of logical error rates, acceptance probabilities, and resource overheads at system sizes beyond the reach of existing methods.

An important direction is to extend the range of supported measurement gadgets.
Some magic state distillation protocols based on gate injection, such as the parity-checking constructions of Campbell and Howard~\cite{campbell2018magic}, fall outside the sufficient conditions of  \cref{thm:iterative-measurements} and therefore require a separate compatibility analysis.
Extending our guarantees to cover these protocols would broaden the scope of the simulator.
Another possibility is to describe states in a rotated local Clifford basis, potentially bringing Hadamard-based cultivation within the framework \cite{chamberland2020very,sahay2026foldtransversal}.
Although such a basis change can map the Hadamard observable to a DCP operator, compatibility must also be established for the transformed preparations, gates, and measurements throughout the protocol.

Finally, the DCP decompositions introduced here suggest a route to circuits that leave the 3PP state class: understanding which fault-tolerant protocols admit small DCP decompositions could extend the practical scope of the simulator beyond the polynomial-time regime established in this work.
